% Options for packages loaded elsewhere
% Options for packages loaded elsewhere
\PassOptionsToPackage{unicode}{hyperref}
\PassOptionsToPackage{hyphens}{url}
\PassOptionsToPackage{dvipsnames,svgnames,x11names}{xcolor}
%
\documentclass[
  letterpaper,
  DIV=11,
  numbers=noendperiod]{scrartcl}
\usepackage{xcolor}
\usepackage{amsmath,amssymb}
\setcounter{secnumdepth}{-\maxdimen} % remove section numbering
\usepackage{iftex}
\ifPDFTeX
  \usepackage[T1]{fontenc}
  \usepackage[utf8]{inputenc}
  \usepackage{textcomp} % provide euro and other symbols
\else % if luatex or xetex
  \usepackage{unicode-math} % this also loads fontspec
  \defaultfontfeatures{Scale=MatchLowercase}
  \defaultfontfeatures[\rmfamily]{Ligatures=TeX,Scale=1}
\fi
\usepackage{lmodern}
\ifPDFTeX\else
  % xetex/luatex font selection
  \setmainfont[]{IBM Plex Sans}
  \setsansfont[]{IBM Plex Sans}
  \setmonofont[]{JuliaMono}
\fi
% Use upquote if available, for straight quotes in verbatim environments
\IfFileExists{upquote.sty}{\usepackage{upquote}}{}
\IfFileExists{microtype.sty}{% use microtype if available
  \usepackage[]{microtype}
  \UseMicrotypeSet[protrusion]{basicmath} % disable protrusion for tt fonts
}{}
\makeatletter
\@ifundefined{KOMAClassName}{% if non-KOMA class
  \IfFileExists{parskip.sty}{%
    \usepackage{parskip}
  }{% else
    \setlength{\parindent}{0pt}
    \setlength{\parskip}{6pt plus 2pt minus 1pt}}
}{% if KOMA class
  \KOMAoptions{parskip=half}}
\makeatother
% Make \paragraph and \subparagraph free-standing
\makeatletter
\ifx\paragraph\undefined\else
  \let\oldparagraph\paragraph
  \renewcommand{\paragraph}{
    \@ifstar
      \xxxParagraphStar
      \xxxParagraphNoStar
  }
  \newcommand{\xxxParagraphStar}[1]{\oldparagraph*{#1}\mbox{}}
  \newcommand{\xxxParagraphNoStar}[1]{\oldparagraph{#1}\mbox{}}
\fi
\ifx\subparagraph\undefined\else
  \let\oldsubparagraph\subparagraph
  \renewcommand{\subparagraph}{
    \@ifstar
      \xxxSubParagraphStar
      \xxxSubParagraphNoStar
  }
  \newcommand{\xxxSubParagraphStar}[1]{\oldsubparagraph*{#1}\mbox{}}
  \newcommand{\xxxSubParagraphNoStar}[1]{\oldsubparagraph{#1}\mbox{}}
\fi
\makeatother

\usepackage{color}
\usepackage{fancyvrb}
\newcommand{\VerbBar}{|}
\newcommand{\VERB}{\Verb[commandchars=\\\{\}]}
\DefineVerbatimEnvironment{Highlighting}{Verbatim}{commandchars=\\\{\}}
% Add ',fontsize=\small' for more characters per line
\usepackage{framed}
\definecolor{shadecolor}{RGB}{241,243,245}
\newenvironment{Shaded}{\begin{snugshade}}{\end{snugshade}}
\newcommand{\AlertTok}[1]{\textcolor[rgb]{0.68,0.00,0.00}{#1}}
\newcommand{\AnnotationTok}[1]{\textcolor[rgb]{0.37,0.37,0.37}{#1}}
\newcommand{\AttributeTok}[1]{\textcolor[rgb]{0.40,0.45,0.13}{#1}}
\newcommand{\BaseNTok}[1]{\textcolor[rgb]{0.68,0.00,0.00}{#1}}
\newcommand{\BuiltInTok}[1]{\textcolor[rgb]{0.00,0.23,0.31}{#1}}
\newcommand{\CharTok}[1]{\textcolor[rgb]{0.13,0.47,0.30}{#1}}
\newcommand{\CommentTok}[1]{\textcolor[rgb]{0.37,0.37,0.37}{#1}}
\newcommand{\CommentVarTok}[1]{\textcolor[rgb]{0.37,0.37,0.37}{\textit{#1}}}
\newcommand{\ConstantTok}[1]{\textcolor[rgb]{0.56,0.35,0.01}{#1}}
\newcommand{\ControlFlowTok}[1]{\textcolor[rgb]{0.00,0.23,0.31}{\textbf{#1}}}
\newcommand{\DataTypeTok}[1]{\textcolor[rgb]{0.68,0.00,0.00}{#1}}
\newcommand{\DecValTok}[1]{\textcolor[rgb]{0.68,0.00,0.00}{#1}}
\newcommand{\DocumentationTok}[1]{\textcolor[rgb]{0.37,0.37,0.37}{\textit{#1}}}
\newcommand{\ErrorTok}[1]{\textcolor[rgb]{0.68,0.00,0.00}{#1}}
\newcommand{\ExtensionTok}[1]{\textcolor[rgb]{0.00,0.23,0.31}{#1}}
\newcommand{\FloatTok}[1]{\textcolor[rgb]{0.68,0.00,0.00}{#1}}
\newcommand{\FunctionTok}[1]{\textcolor[rgb]{0.28,0.35,0.67}{#1}}
\newcommand{\ImportTok}[1]{\textcolor[rgb]{0.00,0.46,0.62}{#1}}
\newcommand{\InformationTok}[1]{\textcolor[rgb]{0.37,0.37,0.37}{#1}}
\newcommand{\KeywordTok}[1]{\textcolor[rgb]{0.00,0.23,0.31}{\textbf{#1}}}
\newcommand{\NormalTok}[1]{\textcolor[rgb]{0.00,0.23,0.31}{#1}}
\newcommand{\OperatorTok}[1]{\textcolor[rgb]{0.37,0.37,0.37}{#1}}
\newcommand{\OtherTok}[1]{\textcolor[rgb]{0.00,0.23,0.31}{#1}}
\newcommand{\PreprocessorTok}[1]{\textcolor[rgb]{0.68,0.00,0.00}{#1}}
\newcommand{\RegionMarkerTok}[1]{\textcolor[rgb]{0.00,0.23,0.31}{#1}}
\newcommand{\SpecialCharTok}[1]{\textcolor[rgb]{0.37,0.37,0.37}{#1}}
\newcommand{\SpecialStringTok}[1]{\textcolor[rgb]{0.13,0.47,0.30}{#1}}
\newcommand{\StringTok}[1]{\textcolor[rgb]{0.13,0.47,0.30}{#1}}
\newcommand{\VariableTok}[1]{\textcolor[rgb]{0.07,0.07,0.07}{#1}}
\newcommand{\VerbatimStringTok}[1]{\textcolor[rgb]{0.13,0.47,0.30}{#1}}
\newcommand{\WarningTok}[1]{\textcolor[rgb]{0.37,0.37,0.37}{\textit{#1}}}

\usepackage{longtable,booktabs,array}
\usepackage{calc} % for calculating minipage widths
% Correct order of tables after \paragraph or \subparagraph
\usepackage{etoolbox}
\makeatletter
\patchcmd\longtable{\par}{\if@noskipsec\mbox{}\fi\par}{}{}
\makeatother
% Allow footnotes in longtable head/foot
\IfFileExists{footnotehyper.sty}{\usepackage{footnotehyper}}{\usepackage{footnote}}
\makesavenoteenv{longtable}
\usepackage{graphicx}
\makeatletter
\newsavebox\pandoc@box
\newcommand*\pandocbounded[1]{% scales image to fit in text height/width
  \sbox\pandoc@box{#1}%
  \Gscale@div\@tempa{\textheight}{\dimexpr\ht\pandoc@box+\dp\pandoc@box\relax}%
  \Gscale@div\@tempb{\linewidth}{\wd\pandoc@box}%
  \ifdim\@tempb\p@<\@tempa\p@\let\@tempa\@tempb\fi% select the smaller of both
  \ifdim\@tempa\p@<\p@\scalebox{\@tempa}{\usebox\pandoc@box}%
  \else\usebox{\pandoc@box}%
  \fi%
}
% Set default figure placement to htbp
\def\fps@figure{htbp}
\makeatother





\setlength{\emergencystretch}{3em} % prevent overfull lines

\providecommand{\tightlist}{%
  \setlength{\itemsep}{0pt}\setlength{\parskip}{0pt}}



 


\usepackage{fvextra}
\DefineVerbatimEnvironment{Highlighting}{Verbatim}{breaklines,commandchars=\\\{\}}
\KOMAoption{captions}{tableheading}
\makeatletter
\@ifpackageloaded{caption}{}{\usepackage{caption}}
\AtBeginDocument{%
\ifdefined\contentsname
  \renewcommand*\contentsname{Table of contents}
\else
  \newcommand\contentsname{Table of contents}
\fi
\ifdefined\listfigurename
  \renewcommand*\listfigurename{List of Figures}
\else
  \newcommand\listfigurename{List of Figures}
\fi
\ifdefined\listtablename
  \renewcommand*\listtablename{List of Tables}
\else
  \newcommand\listtablename{List of Tables}
\fi
\ifdefined\figurename
  \renewcommand*\figurename{Figure}
\else
  \newcommand\figurename{Figure}
\fi
\ifdefined\tablename
  \renewcommand*\tablename{Table}
\else
  \newcommand\tablename{Table}
\fi
}
\@ifpackageloaded{float}{}{\usepackage{float}}
\floatstyle{ruled}
\@ifundefined{c@chapter}{\newfloat{codelisting}{h}{lop}}{\newfloat{codelisting}{h}{lop}[chapter]}
\floatname{codelisting}{Listing}
\newcommand*\listoflistings{\listof{codelisting}{List of Listings}}
\makeatother
\makeatletter
\makeatother
\makeatletter
\@ifpackageloaded{caption}{}{\usepackage{caption}}
\@ifpackageloaded{subcaption}{}{\usepackage{subcaption}}
\makeatother
\makeatletter
\definecolor{QuartoInternalColor2}{rgb}{0.00,0.45,0.15}
\definecolor{QuartoInternalColor1}{rgb}{0,0,0}
\makeatother
\makeatletter
\@ifpackageloaded{tikz}{}{\usepackage{tikz}}
\makeatother
        \newcommand*\circled[1]{\tikz[baseline=(char.base)]{
          \node[shape=circle,draw,inner sep=1pt] (char) {{\scriptsize#1}};}}  
                  
\usepackage{bookmark}
\IfFileExists{xurl.sty}{\usepackage{xurl}}{} % add URL line breaks if available
\urlstyle{same}
\hypersetup{
  pdftitle={Homework 4 Solutions},
  colorlinks=true,
  linkcolor={blue},
  filecolor={Maroon},
  citecolor={Blue},
  urlcolor={Blue},
  pdfcreator={LaTeX via pandoc}}


\title{Homework 4 Solutions}
\usepackage{etoolbox}
\makeatletter
\providecommand{\subtitle}[1]{% add subtitle to \maketitle
  \apptocmd{\@title}{\par {\large #1 \par}}{}{}
}
\makeatother
\subtitle{BEE 4850/5850}
\author{}
\date{}
\begin{document}
\maketitle

\RecustomVerbatimEnvironment{verbatim}{Verbatim}{
showspaces = false,
showtabs = false,
breaksymbolleft={},
breaklines
% Note: setting commandchars=\\\{\} here will cause an error
}


\subsection{Overview}\label{overview}

\subsubsection{Instructions}\label{instructions}

The goal of this homework assignment is to practice model evaluation and
using cross-validation and information criteria to distinguish between
models.

\begin{itemize}
\tightlist
\item
  Problem 1 asks you to fit a model for cherry blossom bloom dates and
  use cross-validation to draw conclusions about generalizability.
\item
  Problem 2 asks you to compare several linear regression models for
  environmental influences on non-accidental death data.
\end{itemize}

\subsubsection{Load Environment}\label{load-environment}

The following code loads the environment and makes sure all needed
packages are installed. This should be at the start of most Julia
scripts.

\begin{Shaded}
\begin{Highlighting}[]
\ImportTok{import} \BuiltInTok{Pkg}
\BuiltInTok{Pkg}\NormalTok{.}\FunctionTok{activate}\NormalTok{(}\PreprocessorTok{@\_\_DIR\_\_}\NormalTok{)}
\BuiltInTok{Pkg}\NormalTok{.}\FunctionTok{instantiate}\NormalTok{()}
\end{Highlighting}
\end{Shaded}

The following packages are included in the environment (to help you find
other similar packages in other languages). The code below loads these
packages for use in the subsequent notebook (the desired functionality
for each package is commented next to the package).

\begin{Shaded}
\begin{Highlighting}[]
\ImportTok{using} \BuiltInTok{Random} \CommentTok{\# random number generation and seed{-}setting}
\ImportTok{using} \BuiltInTok{DataFrames} \CommentTok{\# tabular data structure}
\ImportTok{using} \BuiltInTok{DataFramesMeta} \CommentTok{\# API which can simplify chains of DataFrames transformations}
\ImportTok{using} \BuiltInTok{CSV\#}\NormalTok{ reads/writes .csv files}
\ImportTok{using} \BuiltInTok{Distributions} \CommentTok{\# interface to work with probability distributions}
\ImportTok{using} \BuiltInTok{Plots} \CommentTok{\# plotting library}
\ImportTok{using} \BuiltInTok{StatsBase} \CommentTok{\# statistical quantities like mean, median, etc}
\ImportTok{using} \BuiltInTok{StatsPlots} \CommentTok{\# some additional statistical plotting tools}
\ImportTok{using} \BuiltInTok{Optim}
\ImportTok{using} \BuiltInTok{LaTeXStrings}
\ImportTok{using} \BuiltInTok{Dates}

\BuiltInTok{Random}\NormalTok{.}\FunctionTok{seed!}\NormalTok{(}\FloatTok{1}\NormalTok{)}
\end{Highlighting}
\end{Shaded}

\subsection{Problems}\label{problems}

\subsubsection{Problem 1}\label{problem-1}

\paragraph{Problem 1.1}\label{problem-1.1}

Let's load the data from `data/HadCRUT5.1Analysis\_gl.txt' and
`data/kyoto.csv' and merge the two datasets along common years.

\phantomsection\label{annotated-cell-3}%
\begin{Shaded}
\begin{Highlighting}[]
\NormalTok{temp\_dat }\OperatorTok{=}\NormalTok{ CSV.}\FunctionTok{read}\NormalTok{(}\FunctionTok{joinpath}\NormalTok{(}\StringTok{"data"}\NormalTok{, }\StringTok{"HadCRUT5.1Analysis\_gl.txt"}\NormalTok{), DataFrame, delim}\OperatorTok{=}\StringTok{" "}\NormalTok{, ignorerepeated}\OperatorTok{=}\ConstantTok{true}\NormalTok{, header}\OperatorTok{=}\ConstantTok{false}\NormalTok{) }\hspace*{\fill}\NormalTok{\circled{1}}
\NormalTok{temp\_dat }\OperatorTok{=}\NormalTok{ temp\_dat[}\FloatTok{1}\OperatorTok{:}\FloatTok{2}\OperatorTok{:}\FunctionTok{nrow}\NormalTok{(temp\_dat), [}\FloatTok{1}\NormalTok{, }\FloatTok{4}\NormalTok{]] }\hspace*{\fill}\NormalTok{\circled{2}}
\FunctionTok{rename!}\NormalTok{(temp\_dat, [}\StringTok{"Year"}\NormalTok{, }\StringTok{"March"}\NormalTok{])}

\NormalTok{kyoto\_dat }\OperatorTok{=}\NormalTok{ CSV.}\FunctionTok{read}\NormalTok{(}\FunctionTok{joinpath}\NormalTok{(}\StringTok{"data"}\NormalTok{, }\StringTok{"kyoto.csv"}\NormalTok{), DataFrame, missingstring}\OperatorTok{=}\StringTok{"NA"}\NormalTok{)}
\NormalTok{kyoto\_dat }\OperatorTok{=}\NormalTok{ kyoto\_dat[}\OperatorTok{:}\NormalTok{, }\FloatTok{1}\OperatorTok{:}\FloatTok{2}\NormalTok{]}
\FunctionTok{rename!}\NormalTok{(kyoto\_dat, [}\OperatorTok{:}\NormalTok{Year, }\OperatorTok{:}\NormalTok{DOY]) }\hspace*{\fill}\NormalTok{\circled{3}}

\NormalTok{dat }\OperatorTok{=} \FunctionTok{innerjoin}\NormalTok{(kyoto\_dat, temp\_dat, on }\OperatorTok{=} \OperatorTok{:}\NormalTok{Year) }\hspace*{\fill}\NormalTok{\circled{4}}
\NormalTok{dat  }\OperatorTok{=}\NormalTok{ dat[}\FunctionTok{completecases}\NormalTok{(dat), }\OperatorTok{:}\NormalTok{] }\hspace*{\fill}\NormalTok{\circled{5}}
\FunctionTok{disallowmissing!}\NormalTok{(dat) }
\end{Highlighting}
\end{Shaded}

\begin{description}
\tightlist
\item[\circled{1}]
The data is space-delimited and the number of spaces between columns
differs in the temperature (odd) and number of station (even) rows, so
we use \texttt{ignorerepeated=true} to cause the parser to not interpret
multiple consecutive spaces as delimiting multiple columns.
\item[\circled{2}]
We need to only keep the odd rows (which are the temperatures) and the
year/March anomaly columns.
\item[\circled{3}]
The periods in the column names in \texttt{kyoto.csv} are annoying for
later \texttt{DataFrame} operations, so we'll just rename the columns to
get rid of them.
\item[\circled{4}]
An \texttt{innerjoin} is a type of database join that merges two
datasets based on the entries of a common ID column which are present in
both datasets.
\item[\circled{5}]
There are still some years with \texttt{missing} DOYs, so we use
\texttt{completecases} to remove those and \texttt{disallowmissing!} to
change the column datatype.
\end{description}

\begin{verbatim}
┌ Warning: thread = 1 warning: only found 13 / 14 columns around data row: 2. Filling remaining columns with `missing`
└ @ CSV ~/.julia/packages/CSV/LiiJM/src/file.jl:592
┌ Warning: thread = 1 warning: only found 13 / 14 columns around data row: 4. Filling remaining columns with `missing`
└ @ CSV ~/.julia/packages/CSV/LiiJM/src/file.jl:592
┌ Warning: thread = 1 warning: only found 13 / 14 columns around data row: 6. Filling remaining columns with `missing`
└ @ CSV ~/.julia/packages/CSV/LiiJM/src/file.jl:592
┌ Warning: thread = 1 warning: only found 13 / 14 columns around data row: 8. Filling remaining columns with `missing`
└ @ CSV ~/.julia/packages/CSV/LiiJM/src/file.jl:592
┌ Warning: thread = 1 warning: only found 13 / 14 columns around data row: 10. Filling remaining columns with `missing`
└ @ CSV ~/.julia/packages/CSV/LiiJM/src/file.jl:592
┌ Warning: thread = 1 warning: only found 13 / 14 columns around data row: 12. Filling remaining columns with `missing`
└ @ CSV ~/.julia/packages/CSV/LiiJM/src/file.jl:592
┌ Warning: thread = 1 warning: only found 13 / 14 columns around data row: 14. Filling remaining columns with `missing`
└ @ CSV ~/.julia/packages/CSV/LiiJM/src/file.jl:592
┌ Warning: thread = 1 warning: only found 13 / 14 columns around data row: 16. Filling remaining columns with `missing`
└ @ CSV ~/.julia/packages/CSV/LiiJM/src/file.jl:592
┌ Warning: thread = 1 warning: only found 13 / 14 columns around data row: 18. Filling remaining columns with `missing`
└ @ CSV ~/.julia/packages/CSV/LiiJM/src/file.jl:592
┌ Warning: thread = 1 warning: only found 13 / 14 columns around data row: 20. Filling remaining columns with `missing`
└ @ CSV ~/.julia/packages/CSV/LiiJM/src/file.jl:592
┌ Warning: thread = 1 warning: only found 13 / 14 columns around data row: 22. Filling remaining columns with `missing`
└ @ CSV ~/.julia/packages/CSV/LiiJM/src/file.jl:592
┌ Warning: thread = 1 warning: only found 13 / 14 columns around data row: 24. Filling remaining columns with `missing`
└ @ CSV ~/.julia/packages/CSV/LiiJM/src/file.jl:592
┌ Warning: thread = 1 warning: only found 13 / 14 columns around data row: 26. Filling remaining columns with `missing`
└ @ CSV ~/.julia/packages/CSV/LiiJM/src/file.jl:592
┌ Warning: thread = 1 warning: only found 13 / 14 columns around data row: 28. Filling remaining columns with `missing`
└ @ CSV ~/.julia/packages/CSV/LiiJM/src/file.jl:592
┌ Warning: thread = 1 warning: only found 13 / 14 columns around data row: 30. Filling remaining columns with `missing`
└ @ CSV ~/.julia/packages/CSV/LiiJM/src/file.jl:592
┌ Warning: thread = 1 warning: only found 13 / 14 columns around data row: 32. Filling remaining columns with `missing`
└ @ CSV ~/.julia/packages/CSV/LiiJM/src/file.jl:592
┌ Warning: thread = 1 warning: only found 13 / 14 columns around data row: 34. Filling remaining columns with `missing`
└ @ CSV ~/.julia/packages/CSV/LiiJM/src/file.jl:592
┌ Warning: thread = 1 warning: only found 13 / 14 columns around data row: 36. Filling remaining columns with `missing`
└ @ CSV ~/.julia/packages/CSV/LiiJM/src/file.jl:592
┌ Warning: thread = 1 warning: only found 13 / 14 columns around data row: 38. Filling remaining columns with `missing`
└ @ CSV ~/.julia/packages/CSV/LiiJM/src/file.jl:592
┌ Warning: thread = 1 warning: only found 13 / 14 columns around data row: 40. Filling remaining columns with `missing`
└ @ CSV ~/.julia/packages/CSV/LiiJM/src/file.jl:592
┌ Warning: thread = 1 warning: only found 13 / 14 columns around data row: 42. Filling remaining columns with `missing`
└ @ CSV ~/.julia/packages/CSV/LiiJM/src/file.jl:592
┌ Warning: thread = 1 warning: only found 13 / 14 columns around data row: 44. Filling remaining columns with `missing`
└ @ CSV ~/.julia/packages/CSV/LiiJM/src/file.jl:592
┌ Warning: thread = 1 warning: only found 13 / 14 columns around data row: 46. Filling remaining columns with `missing`
└ @ CSV ~/.julia/packages/CSV/LiiJM/src/file.jl:592
┌ Warning: thread = 1 warning: only found 13 / 14 columns around data row: 48. Filling remaining columns with `missing`
└ @ CSV ~/.julia/packages/CSV/LiiJM/src/file.jl:592
┌ Warning: thread = 1 warning: only found 13 / 14 columns around data row: 50. Filling remaining columns with `missing`
└ @ CSV ~/.julia/packages/CSV/LiiJM/src/file.jl:592
┌ Warning: thread = 1 warning: only found 13 / 14 columns around data row: 52. Filling remaining columns with `missing`
└ @ CSV ~/.julia/packages/CSV/LiiJM/src/file.jl:592
┌ Warning: thread = 1 warning: only found 13 / 14 columns around data row: 54. Filling remaining columns with `missing`
└ @ CSV ~/.julia/packages/CSV/LiiJM/src/file.jl:592
┌ Warning: thread = 1 warning: only found 13 / 14 columns around data row: 56. Filling remaining columns with `missing`
└ @ CSV ~/.julia/packages/CSV/LiiJM/src/file.jl:592
┌ Warning: thread = 1 warning: only found 13 / 14 columns around data row: 58. Filling remaining columns with `missing`
└ @ CSV ~/.julia/packages/CSV/LiiJM/src/file.jl:592
┌ Warning: thread = 1 warning: only found 13 / 14 columns around data row: 60. Filling remaining columns with `missing`
└ @ CSV ~/.julia/packages/CSV/LiiJM/src/file.jl:592
┌ Warning: thread = 1 warning: only found 13 / 14 columns around data row: 62. Filling remaining columns with `missing`
└ @ CSV ~/.julia/packages/CSV/LiiJM/src/file.jl:592
┌ Warning: thread = 1 warning: only found 13 / 14 columns around data row: 64. Filling remaining columns with `missing`
└ @ CSV ~/.julia/packages/CSV/LiiJM/src/file.jl:592
┌ Warning: thread = 1 warning: only found 13 / 14 columns around data row: 66. Filling remaining columns with `missing`
└ @ CSV ~/.julia/packages/CSV/LiiJM/src/file.jl:592
┌ Warning: thread = 1 warning: only found 13 / 14 columns around data row: 68. Filling remaining columns with `missing`
└ @ CSV ~/.julia/packages/CSV/LiiJM/src/file.jl:592
┌ Warning: thread = 1 warning: only found 13 / 14 columns around data row: 70. Filling remaining columns with `missing`
└ @ CSV ~/.julia/packages/CSV/LiiJM/src/file.jl:592
┌ Warning: thread = 1 warning: only found 13 / 14 columns around data row: 72. Filling remaining columns with `missing`
└ @ CSV ~/.julia/packages/CSV/LiiJM/src/file.jl:592
┌ Warning: thread = 1 warning: only found 13 / 14 columns around data row: 74. Filling remaining columns with `missing`
└ @ CSV ~/.julia/packages/CSV/LiiJM/src/file.jl:592
┌ Warning: thread = 1 warning: only found 13 / 14 columns around data row: 76. Filling remaining columns with `missing`
└ @ CSV ~/.julia/packages/CSV/LiiJM/src/file.jl:592
┌ Warning: thread = 1 warning: only found 13 / 14 columns around data row: 78. Filling remaining columns with `missing`
└ @ CSV ~/.julia/packages/CSV/LiiJM/src/file.jl:592
┌ Warning: thread = 1 warning: only found 13 / 14 columns around data row: 80. Filling remaining columns with `missing`
└ @ CSV ~/.julia/packages/CSV/LiiJM/src/file.jl:592
┌ Warning: thread = 1 warning: only found 13 / 14 columns around data row: 82. Filling remaining columns with `missing`
└ @ CSV ~/.julia/packages/CSV/LiiJM/src/file.jl:592
┌ Warning: thread = 1 warning: only found 13 / 14 columns around data row: 84. Filling remaining columns with `missing`
└ @ CSV ~/.julia/packages/CSV/LiiJM/src/file.jl:592
┌ Warning: thread = 1 warning: only found 13 / 14 columns around data row: 86. Filling remaining columns with `missing`
└ @ CSV ~/.julia/packages/CSV/LiiJM/src/file.jl:592
┌ Warning: thread = 1 warning: only found 13 / 14 columns around data row: 88. Filling remaining columns with `missing`
└ @ CSV ~/.julia/packages/CSV/LiiJM/src/file.jl:592
┌ Warning: thread = 1 warning: only found 13 / 14 columns around data row: 90. Filling remaining columns with `missing`
└ @ CSV ~/.julia/packages/CSV/LiiJM/src/file.jl:592
┌ Warning: thread = 1 warning: only found 13 / 14 columns around data row: 92. Filling remaining columns with `missing`
└ @ CSV ~/.julia/packages/CSV/LiiJM/src/file.jl:592
┌ Warning: thread = 1 warning: only found 13 / 14 columns around data row: 94. Filling remaining columns with `missing`
└ @ CSV ~/.julia/packages/CSV/LiiJM/src/file.jl:592
┌ Warning: thread = 1 warning: only found 13 / 14 columns around data row: 96. Filling remaining columns with `missing`
└ @ CSV ~/.julia/packages/CSV/LiiJM/src/file.jl:592
┌ Warning: thread = 1 warning: only found 13 / 14 columns around data row: 98. Filling remaining columns with `missing`
└ @ CSV ~/.julia/packages/CSV/LiiJM/src/file.jl:592
┌ Warning: thread = 1 warning: only found 13 / 14 columns around data row: 100. Filling remaining columns with `missing`
└ @ CSV ~/.julia/packages/CSV/LiiJM/src/file.jl:592
┌ Warning: thread = 1 warning: only found 13 / 14 columns around data row: 102. Filling remaining columns with `missing`
└ @ CSV ~/.julia/packages/CSV/LiiJM/src/file.jl:592
┌ Warning: thread = 1 warning: only found 13 / 14 columns around data row: 104. Filling remaining columns with `missing`
└ @ CSV ~/.julia/packages/CSV/LiiJM/src/file.jl:592
┌ Warning: thread = 1 warning: only found 13 / 14 columns around data row: 106. Filling remaining columns with `missing`
└ @ CSV ~/.julia/packages/CSV/LiiJM/src/file.jl:592
┌ Warning: thread = 1 warning: only found 13 / 14 columns around data row: 108. Filling remaining columns with `missing`
└ @ CSV ~/.julia/packages/CSV/LiiJM/src/file.jl:592
┌ Warning: thread = 1 warning: only found 13 / 14 columns around data row: 110. Filling remaining columns with `missing`
└ @ CSV ~/.julia/packages/CSV/LiiJM/src/file.jl:592
┌ Warning: thread = 1 warning: only found 13 / 14 columns around data row: 112. Filling remaining columns with `missing`
└ @ CSV ~/.julia/packages/CSV/LiiJM/src/file.jl:592
┌ Warning: thread = 1 warning: only found 13 / 14 columns around data row: 114. Filling remaining columns with `missing`
└ @ CSV ~/.julia/packages/CSV/LiiJM/src/file.jl:592
┌ Warning: thread = 1 warning: only found 13 / 14 columns around data row: 116. Filling remaining columns with `missing`
└ @ CSV ~/.julia/packages/CSV/LiiJM/src/file.jl:592
┌ Warning: thread = 1 warning: only found 13 / 14 columns around data row: 118. Filling remaining columns with `missing`
└ @ CSV ~/.julia/packages/CSV/LiiJM/src/file.jl:592
┌ Warning: thread = 1 warning: only found 13 / 14 columns around data row: 120. Filling remaining columns with `missing`
└ @ CSV ~/.julia/packages/CSV/LiiJM/src/file.jl:592
┌ Warning: thread = 1 warning: only found 13 / 14 columns around data row: 122. Filling remaining columns with `missing`
└ @ CSV ~/.julia/packages/CSV/LiiJM/src/file.jl:592
┌ Warning: thread = 1 warning: only found 13 / 14 columns around data row: 124. Filling remaining columns with `missing`
└ @ CSV ~/.julia/packages/CSV/LiiJM/src/file.jl:592
┌ Warning: thread = 1 warning: only found 13 / 14 columns around data row: 126. Filling remaining columns with `missing`
└ @ CSV ~/.julia/packages/CSV/LiiJM/src/file.jl:592
┌ Warning: thread = 1 warning: only found 13 / 14 columns around data row: 128. Filling remaining columns with `missing`
└ @ CSV ~/.julia/packages/CSV/LiiJM/src/file.jl:592
┌ Warning: thread = 1 warning: only found 13 / 14 columns around data row: 130. Filling remaining columns with `missing`
└ @ CSV ~/.julia/packages/CSV/LiiJM/src/file.jl:592
┌ Warning: thread = 1 warning: only found 13 / 14 columns around data row: 132. Filling remaining columns with `missing`
└ @ CSV ~/.julia/packages/CSV/LiiJM/src/file.jl:592
┌ Warning: thread = 1 warning: only found 13 / 14 columns around data row: 134. Filling remaining columns with `missing`
└ @ CSV ~/.julia/packages/CSV/LiiJM/src/file.jl:592
┌ Warning: thread = 1 warning: only found 13 / 14 columns around data row: 136. Filling remaining columns with `missing`
└ @ CSV ~/.julia/packages/CSV/LiiJM/src/file.jl:592
┌ Warning: thread = 1 warning: only found 13 / 14 columns around data row: 138. Filling remaining columns with `missing`
└ @ CSV ~/.julia/packages/CSV/LiiJM/src/file.jl:592
┌ Warning: thread = 1 warning: only found 13 / 14 columns around data row: 140. Filling remaining columns with `missing`
└ @ CSV ~/.julia/packages/CSV/LiiJM/src/file.jl:592
┌ Warning: thread = 1 warning: only found 13 / 14 columns around data row: 142. Filling remaining columns with `missing`
└ @ CSV ~/.julia/packages/CSV/LiiJM/src/file.jl:592
┌ Warning: thread = 1 warning: only found 13 / 14 columns around data row: 144. Filling remaining columns with `missing`
└ @ CSV ~/.julia/packages/CSV/LiiJM/src/file.jl:592
┌ Warning: thread = 1 warning: only found 13 / 14 columns around data row: 146. Filling remaining columns with `missing`
└ @ CSV ~/.julia/packages/CSV/LiiJM/src/file.jl:592
┌ Warning: thread = 1 warning: only found 13 / 14 columns around data row: 148. Filling remaining columns with `missing`
└ @ CSV ~/.julia/packages/CSV/LiiJM/src/file.jl:592
┌ Warning: thread = 1 warning: only found 13 / 14 columns around data row: 150. Filling remaining columns with `missing`
└ @ CSV ~/.julia/packages/CSV/LiiJM/src/file.jl:592
┌ Warning: thread = 1 warning: only found 13 / 14 columns around data row: 152. Filling remaining columns with `missing`
└ @ CSV ~/.julia/packages/CSV/LiiJM/src/file.jl:592
┌ Warning: thread = 1 warning: only found 13 / 14 columns around data row: 154. Filling remaining columns with `missing`
└ @ CSV ~/.julia/packages/CSV/LiiJM/src/file.jl:592
┌ Warning: thread = 1 warning: only found 13 / 14 columns around data row: 156. Filling remaining columns with `missing`
└ @ CSV ~/.julia/packages/CSV/LiiJM/src/file.jl:592
┌ Warning: thread = 1 warning: only found 13 / 14 columns around data row: 158. Filling remaining columns with `missing`
└ @ CSV ~/.julia/packages/CSV/LiiJM/src/file.jl:592
┌ Warning: thread = 1 warning: only found 13 / 14 columns around data row: 160. Filling remaining columns with `missing`
└ @ CSV ~/.julia/packages/CSV/LiiJM/src/file.jl:592
┌ Warning: thread = 1 warning: only found 13 / 14 columns around data row: 162. Filling remaining columns with `missing`
└ @ CSV ~/.julia/packages/CSV/LiiJM/src/file.jl:592
┌ Warning: thread = 1 warning: only found 13 / 14 columns around data row: 164. Filling remaining columns with `missing`
└ @ CSV ~/.julia/packages/CSV/LiiJM/src/file.jl:592
┌ Warning: thread = 1 warning: only found 13 / 14 columns around data row: 166. Filling remaining columns with `missing`
└ @ CSV ~/.julia/packages/CSV/LiiJM/src/file.jl:592
┌ Warning: thread = 1 warning: only found 13 / 14 columns around data row: 168. Filling remaining columns with `missing`
└ @ CSV ~/.julia/packages/CSV/LiiJM/src/file.jl:592
┌ Warning: thread = 1 warning: only found 13 / 14 columns around data row: 170. Filling remaining columns with `missing`
└ @ CSV ~/.julia/packages/CSV/LiiJM/src/file.jl:592
┌ Warning: thread = 1 warning: only found 13 / 14 columns around data row: 172. Filling remaining columns with `missing`
└ @ CSV ~/.julia/packages/CSV/LiiJM/src/file.jl:592
┌ Warning: thread = 1 warning: only found 13 / 14 columns around data row: 174. Filling remaining columns with `missing`
└ @ CSV ~/.julia/packages/CSV/LiiJM/src/file.jl:592
┌ Warning: thread = 1 warning: only found 13 / 14 columns around data row: 176. Filling remaining columns with `missing`
└ @ CSV ~/.julia/packages/CSV/LiiJM/src/file.jl:592
┌ Warning: thread = 1 warning: only found 13 / 14 columns around data row: 178. Filling remaining columns with `missing`
└ @ CSV ~/.julia/packages/CSV/LiiJM/src/file.jl:592
┌ Warning: thread = 1 warning: only found 13 / 14 columns around data row: 180. Filling remaining columns with `missing`
└ @ CSV ~/.julia/packages/CSV/LiiJM/src/file.jl:592
┌ Warning: thread = 1 warning: only found 13 / 14 columns around data row: 182. Filling remaining columns with `missing`
└ @ CSV ~/.julia/packages/CSV/LiiJM/src/file.jl:592
┌ Warning: thread = 1 warning: only found 13 / 14 columns around data row: 184. Filling remaining columns with `missing`
└ @ CSV ~/.julia/packages/CSV/LiiJM/src/file.jl:592
┌ Warning: thread = 1 warning: only found 13 / 14 columns around data row: 186. Filling remaining columns with `missing`
└ @ CSV ~/.julia/packages/CSV/LiiJM/src/file.jl:592
┌ Warning: thread = 1 warning: only found 13 / 14 columns around data row: 188. Filling remaining columns with `missing`
└ @ CSV ~/.julia/packages/CSV/LiiJM/src/file.jl:592
┌ Warning: thread = 1 warning: only found 13 / 14 columns around data row: 190. Filling remaining columns with `missing`
└ @ CSV ~/.julia/packages/CSV/LiiJM/src/file.jl:592
┌ Warning: thread = 1 warning: only found 13 / 14 columns around data row: 192. Filling remaining columns with `missing`
└ @ CSV ~/.julia/packages/CSV/LiiJM/src/file.jl:592
┌ Warning: thread = 1 warning: only found 13 / 14 columns around data row: 194. Filling remaining columns with `missing`
└ @ CSV ~/.julia/packages/CSV/LiiJM/src/file.jl:592
┌ Warning: thread = 1 warning: only found 13 / 14 columns around data row: 196. Filling remaining columns with `missing`
└ @ CSV ~/.julia/packages/CSV/LiiJM/src/file.jl:592
┌ Warning: thread = 1 warning: only found 13 / 14 columns around data row: 198. Filling remaining columns with `missing`
└ @ CSV ~/.julia/packages/CSV/LiiJM/src/file.jl:592
┌ Warning: thread = 1 warning: only found 13 / 14 columns around data row: 200. Filling remaining columns with `missing`
└ @ CSV ~/.julia/packages/CSV/LiiJM/src/file.jl:592
┌ Warning: thread = 1 warning: only found 13 / 14 columns around data row: 202. Filling remaining columns with `missing`
└ @ CSV ~/.julia/packages/CSV/LiiJM/src/file.jl:592
┌ Warning: thread = 1: too many warnings, silencing any further warnings
└ @ CSV ~/.julia/packages/CSV/LiiJM/src/file.jl:597
\end{verbatim}

\begin{tabular}{r|ccc}
    & Year & DOY & March\\
    \hline
    & Int64 & Int64 & Float64\\
    \hline
    1 & 1850 & 109 & -0.627 \\
    2 & 1851 & 102 & -0.641 \\
    3 & 1852 & 113 & -0.567 \\
    4 & 1853 & 102 & -0.273 \\
    5 & 1854 & 102 & -0.247 \\
    6 & 1855 & 101 & -0.345 \\
    7 & 1856 & 111 & -0.371 \\
    8 & 1857 & 108 & -0.488 \\
    9 & 1858 & 104 & -0.566 \\
    10 & 1859 & 110 & -0.322 \\
    11 & 1860 & 113 & -0.789 \\
    12 & 1861 & 100 & -0.422 \\
    13 & 1862 & 117 & -0.538 \\
    14 & 1863 & 118 & -0.414 \\
    15 & 1864 & 105 & -0.531 \\
    16 & 1865 & 97 & -0.717 \\
    17 & 1866 & 102 & -0.748 \\
    18 & 1867 & 101 & -0.712 \\
    19 & 1868 & 104 & -0.332 \\
    20 & 1869 & 105 & -0.651 \\
    21 & 1870 & 99 & -0.495 \\
    22 & 1871 & 104 & -0.171 \\
    23 & 1873 & 104 & -0.39 \\
    24 & 1874 & 107 & -0.664 \\
    $\dots$ & $\dots$ & $\dots$ & $\dots$ \\
\end{tabular}

Now let's fit the model.

\begin{Shaded}
\begin{Highlighting}[]
\KeywordTok{function} \FunctionTok{cherry\_loglik}\NormalTok{(p, doy, temp)}
\NormalTok{    b0, b1, s }\OperatorTok{=}\NormalTok{ p}
\NormalTok{    m }\OperatorTok{=}\NormalTok{ b0 }\OperatorTok{.+}\NormalTok{ b1 }\OperatorTok{*}\NormalTok{ temp}
\NormalTok{    ll }\OperatorTok{=} \FunctionTok{sum}\NormalTok{(}\FunctionTok{logpdf}\NormalTok{.(}\FunctionTok{Normal}\NormalTok{.(m, s), doy))}
    \ControlFlowTok{return}\NormalTok{ ll}
\KeywordTok{end}

\NormalTok{lb }\OperatorTok{=}\NormalTok{ [}\OperatorTok{{-}}\FloatTok{200}\NormalTok{., }\OperatorTok{{-}}\FloatTok{100}\NormalTok{., }\FloatTok{0.01}\NormalTok{]}
\NormalTok{ub }\OperatorTok{=}\NormalTok{ [}\FloatTok{200}\NormalTok{., }\FloatTok{100}\NormalTok{., }\FloatTok{10}\NormalTok{.]}
\NormalTok{p0 }\OperatorTok{=}\NormalTok{ [}\FloatTok{100}\NormalTok{., }\FloatTok{5}\NormalTok{., }\FloatTok{5}\NormalTok{.]}
\NormalTok{optim\_out }\OperatorTok{=} \FunctionTok{optimize}\NormalTok{(p }\OperatorTok{{-}\textgreater{}} \FunctionTok{{-}cherry\_loglik}\NormalTok{(p, dat.DOY, dat.March), lb, ub, p0)}
\NormalTok{p\_mle }\OperatorTok{=}\NormalTok{ optim\_out.minimizer}
\PreprocessorTok{@show}\NormalTok{ p\_mle;}
\end{Highlighting}
\end{Shaded}

\begin{verbatim}
p_mle = [100.6710269629213, -7.84224390609888, 4.717152358910784]
\end{verbatim}

This suggests that every degree of increased average March temperature
makes the cherry blossoms flower around 8 days earlier.

\paragraph{Problem 1.2}\label{problem-1.2}

We can get the log-loss metric directly from the optimization, as the
negative log-likelihood was the optimization objective.

\begin{Shaded}
\begin{Highlighting}[]
\NormalTok{ll }\OperatorTok{=}\NormalTok{ optim\_out.minimum}
\PreprocessorTok{@show}\NormalTok{ ll;}
\end{Highlighting}
\end{Shaded}

\begin{verbatim}
ll = 478.19315773357795
\end{verbatim}

For the mean-square error, we need to compute the predictions.

\begin{Shaded}
\begin{Highlighting}[]
\NormalTok{pred }\OperatorTok{=}\NormalTok{ p\_mle[}\FloatTok{1}\NormalTok{] }\OperatorTok{.+}\NormalTok{ p\_mle[}\FloatTok{2}\NormalTok{] }\OperatorTok{*}\NormalTok{ dat.March}
\NormalTok{mse }\OperatorTok{=} \FunctionTok{mean}\NormalTok{((pred }\OperatorTok{.{-}}\NormalTok{ dat.DOY)}\OperatorTok{.\^{}}\FloatTok{2}\NormalTok{)}
\PreprocessorTok{@show}\NormalTok{ mse;}
\end{Highlighting}
\end{Shaded}

\begin{verbatim}
mse = 22.251526373148707
\end{verbatim}

\paragraph{Problem 1.3}\label{problem-1.3}

To conduct a 5-fold cross-validation, we need to split the data and then
repeat the earlier workflow and store the metrics. Since we're modeling
this using a linear regression, which assumes that the observations are
independent, we can do an independent split by random partitioning of
the data rows. Let's write a function to do this.

\phantomsection\label{annotated-cell-7}%
\begin{Shaded}
\begin{Highlighting}[]
\KeywordTok{function} \FunctionTok{cherry\_cv}\NormalTok{(dat, nfolds }\OperatorTok{=} \FloatTok{5}\NormalTok{)}
    \CommentTok{\# split the data}
\NormalTok{    labels }\OperatorTok{=} \FunctionTok{repeat}\NormalTok{(}\FloatTok{1}\OperatorTok{:}\NormalTok{nfolds, outer }\OperatorTok{=} \FunctionTok{Int}\NormalTok{(}\FunctionTok{floor}\NormalTok{(}\FunctionTok{nrow}\NormalTok{(dat) }\OperatorTok{/}\NormalTok{ nfolds))) }\hspace*{\fill}\NormalTok{\circled{1}}
\NormalTok{    labels }\OperatorTok{=}\NormalTok{ [labels; }\FloatTok{1}\NormalTok{] }
\NormalTok{    fold\_labels }\OperatorTok{=} \FunctionTok{shuffle}\NormalTok{(labels) }\CommentTok{\# randomly permite the labels}
    \CommentTok{\# create storage for error metrics}
\NormalTok{    cv\_ll }\OperatorTok{=} \FunctionTok{zeros}\NormalTok{(nfolds)}
\NormalTok{    cv\_mse }\OperatorTok{=} \FunctionTok{zeros}\NormalTok{(nfolds)}
    \CommentTok{\# set optimization parameters}
\NormalTok{    lb }\OperatorTok{=}\NormalTok{ [}\OperatorTok{{-}}\FloatTok{200}\NormalTok{., }\OperatorTok{{-}}\FloatTok{100}\NormalTok{., }\FloatTok{0.01}\NormalTok{]}
\NormalTok{    ub }\OperatorTok{=}\NormalTok{ [}\FloatTok{200}\NormalTok{., }\FloatTok{100}\NormalTok{., }\FloatTok{10}\NormalTok{.]}
\NormalTok{    p0 }\OperatorTok{=}\NormalTok{ [}\FloatTok{100}\NormalTok{., }\FloatTok{5}\NormalTok{., }\FloatTok{5}\NormalTok{.]}
    \ControlFlowTok{for}\NormalTok{ fold }\KeywordTok{in} \FloatTok{1}\OperatorTok{:}\NormalTok{nfolds}
\NormalTok{        test\_idx }\OperatorTok{=} \FunctionTok{findall}\NormalTok{(fold\_labels }\OperatorTok{.==}\NormalTok{ fold)}
\NormalTok{        train\_idx }\OperatorTok{=} \FunctionTok{setdiff}\NormalTok{(}\FloatTok{1}\OperatorTok{:}\FunctionTok{nrow}\NormalTok{(dat), test\_idx)}
        \CommentTok{\# fit model to the fold}
\NormalTok{        optim\_out }\OperatorTok{=} \FunctionTok{optimize}\NormalTok{(p }\OperatorTok{{-}\textgreater{}} \FunctionTok{{-}cherry\_loglik}\NormalTok{(p, dat[train\_idx, }\OperatorTok{:}\NormalTok{DOY], dat[train\_idx, }\OperatorTok{:}\NormalTok{March]), lb, ub, p0)}
\NormalTok{        p\_mle }\OperatorTok{=}\NormalTok{ optim\_out.minimizer}
\NormalTok{        pred }\OperatorTok{=}\NormalTok{ p\_mle[}\FloatTok{1}\NormalTok{] }\OperatorTok{.+}\NormalTok{ p\_mle[}\FloatTok{2}\NormalTok{] }\OperatorTok{*}\NormalTok{ dat[test\_idx, }\OperatorTok{:}\NormalTok{March]}
\NormalTok{        cv\_ll[fold] }\OperatorTok{=} \FunctionTok{{-}mean}\NormalTok{(}\FunctionTok{logpdf}\NormalTok{.(}\FunctionTok{Normal}\NormalTok{.(pred, p\_mle[}\FloatTok{3}\NormalTok{]), dat[test\_idx, }\OperatorTok{:}\NormalTok{DOY])) }\hspace*{\fill}\NormalTok{\circled{2}}
\NormalTok{        cv\_mse[fold] }\OperatorTok{=} \FunctionTok{mean}\NormalTok{((pred }\OperatorTok{.{-}}\NormalTok{ dat[test\_idx, }\OperatorTok{:}\NormalTok{DOY])}\OperatorTok{.\^{}}\FloatTok{2}\NormalTok{)}
    \ControlFlowTok{end}
    \ControlFlowTok{return}\NormalTok{ cv\_ll, cv\_mse}
\KeywordTok{end}
\NormalTok{cv\_ll, cv\_mse }\OperatorTok{=} \FunctionTok{cherry\_cv}\NormalTok{(dat, }\FloatTok{5}\NormalTok{)}
\end{Highlighting}
\end{Shaded}

\begin{description}
\tightlist
\item[\circled{1}]
The number of data points is not divisible evenly by 5, so we just
create as close to an even set of splits as we can and then add the last
data point to one of the remaining folds (in this case the first one).
\item[\circled{2}]
We get the mean log score here because the number of data points differs
between the CV folds and the full-data.
\end{description}

\begin{verbatim}
([3.0150202460994526, 3.0283992864632503, 2.8857968018404456, 3.0286095541526157, 2.9459127505005456], [24.20962065059353, 24.765915255574352, 18.289382346763833, 24.772091567081112, 21.155985453526426])
\end{verbatim}

Now let's compare to the total-data/in-sample tests.

\begin{Shaded}
\begin{Highlighting}[]
\PreprocessorTok{@show}\NormalTok{ ll }\OperatorTok{/} \FunctionTok{nrow}\NormalTok{(dat);}
\PreprocessorTok{@show} \FunctionTok{mean}\NormalTok{(cv\_ll);}
\end{Highlighting}
\end{Shaded}

\begin{verbatim}
ll / nrow(dat) = 2.9701438368545214
mean(cv_ll) = 2.980747727811262
\end{verbatim}

\begin{Shaded}
\begin{Highlighting}[]
\PreprocessorTok{@show}\NormalTok{ mse;}
\PreprocessorTok{@show} \FunctionTok{mean}\NormalTok{(cv\_mse);}
\end{Highlighting}
\end{Shaded}

\begin{verbatim}
mse = 22.251526373148707
mean(cv_mse) = 22.63859905470785
\end{verbatim}

We can see that in general, the model's in-sample performance is not too
different than its average CV performance, suggesting that there is no
systematic overfitting occurring.

\paragraph{Problem 1.4}\label{problem-1.4}

Since there is no evidence of overfitting from the cross-validation
test, we might conclude that the model generalizes well.

\subsubsection{Problem 2}\label{problem-2}

\paragraph{Problem 2.1}\label{problem-2.1}

Loading and plotting the deaths data from \texttt{data/chicago.csv}:

\begin{figure}

\centering{

\begin{Shaded}
\begin{Highlighting}[]
\NormalTok{chicago\_dat }\OperatorTok{=}\NormalTok{ CSV.}\FunctionTok{read}\NormalTok{(}\FunctionTok{joinpath}\NormalTok{(}\StringTok{"data"}\NormalTok{, }\StringTok{"chicago.csv"}\NormalTok{), DataFrame, delim}\OperatorTok{=}\StringTok{","}\NormalTok{, missingstring}\OperatorTok{=}\StringTok{"NA"}\NormalTok{, header}\OperatorTok{=}\ConstantTok{true}\NormalTok{)}
\NormalTok{day\_zero }\OperatorTok{=} \FunctionTok{Date}\NormalTok{(}\StringTok{"1993{-}12{-}31"}\NormalTok{)}
\NormalTok{chicago\_dat.}\DataTypeTok{Date} \OperatorTok{=}\NormalTok{ day\_zero }\OperatorTok{.+} \FunctionTok{Day}\NormalTok{.(chicago\_dat.time }\OperatorTok{.+} \FloatTok{0.5}\NormalTok{)}
\end{Highlighting}
\end{Shaded}

\begin{Highlighting}
\textcolor{black}{UndefVarError: UndefVarError(:Date, Main.Notebook)}
\textcolor{black}{UndefVarError: `Date` not defined in `Main.Notebook`}
\textcolor{black}{Suggestion: check for spelling errors or missing imports.}
\textcolor{black}{Hint: a global variable of this name may be made accessible by importing Dates in the current active module Main}
\textcolor{black}{Stacktrace:}
\textcolor{black}{ [1] top-level scope}
\textcolor{black}{}\textcolor{QuartoInternalColor1}{   @}\textcolor{QuartoInternalColor1}{ }\textcolor{QuartoInternalColor1}{\textasciitilde{}/Teaching/BEE4850/solutions/hw04/}\textcolor{QuartoInternalColor1}{}\textcolor{QuartoInternalColor1}{}\textcolor{QuartoInternalColor1}{hw04.qmd:216}\textcolor{QuartoInternalColor1}{}\textcolor{QuartoInternalColor1}{}
\end{Highlighting}

}

\caption{\label{fig-death-data}}

\end{figure}%

To ensure that we use the same observations for the purposes of model
comparison, let's remove the cases corresponding to missing PM2.5 data.

\begin{figure}

\centering{

\begin{Shaded}
\begin{Highlighting}[]
\NormalTok{chicago\_dat }\OperatorTok{=}\NormalTok{ chicago\_dat[.!(ismissing).(chicago\_dat.pm25median), }\OperatorTok{:}\NormalTok{]}

\NormalTok{p\_deaths }\OperatorTok{=} \FunctionTok{plot}\NormalTok{(chicago\_dat.}\DataTypeTok{Date}\NormalTok{, chicago\_dat.death, lw}\OperatorTok{=}\FloatTok{2}\NormalTok{, xlabel}\OperatorTok{=}\StringTok{"Date"}\NormalTok{, ylabel}\OperatorTok{=}\StringTok{"Non{-}Accidental Deaths)"}\NormalTok{, label}\OperatorTok{=}\StringTok{"Data"}\NormalTok{, color}\OperatorTok{=:}\NormalTok{gray)}
\end{Highlighting}
\end{Shaded}

\begin{Highlighting}
\textcolor{black}{ArgumentError: ArgumentError("column name \textbackslash{}"Date\textbackslash{}" not found in the data frame; existing most similar names are: \textbackslash{}"death\textbackslash{}"")}
\textcolor{black}{ArgumentError: column name "Date" not found in the data frame; existing most similar names are: "death"}
\textcolor{black}{Stacktrace:}
\textcolor{black}{ [1] }\textcolor{QuartoInternalColor1}{}\textcolor{QuartoInternalColor1}{lookupname}\textcolor{QuartoInternalColor1}{}
\textcolor{QuartoInternalColor1}{}\textcolor{QuartoInternalColor1}{   @}\textcolor{QuartoInternalColor1}{ }\textcolor{QuartoInternalColor1}{\textasciitilde{}/.julia/packages/DataFrames/b4w9K/src/other/}\textcolor{QuartoInternalColor1}{}\textcolor{QuartoInternalColor1}{}\textcolor{QuartoInternalColor1}{index.jl:434}\textcolor{QuartoInternalColor1}{}\textcolor{QuartoInternalColor1}{}\textcolor{QuartoInternalColor1}{ [inlined]}\textcolor{QuartoInternalColor1}{}
\textcolor{QuartoInternalColor1}{ [2] }\textcolor{QuartoInternalColor1}{}\textcolor{QuartoInternalColor1}{getindex}\textcolor{QuartoInternalColor1}{}
\textcolor{QuartoInternalColor1}{}\textcolor{QuartoInternalColor1}{   @}\textcolor{QuartoInternalColor1}{ }\textcolor{QuartoInternalColor1}{\textasciitilde{}/.julia/packages/DataFrames/b4w9K/src/other/}\textcolor{QuartoInternalColor1}{}\textcolor{QuartoInternalColor1}{}\textcolor{QuartoInternalColor1}{index.jl:440}\textcolor{QuartoInternalColor1}{}\textcolor{QuartoInternalColor1}{}\textcolor{QuartoInternalColor1}{ [inlined]}\textcolor{QuartoInternalColor1}{}
\textcolor{QuartoInternalColor1}{ [3] }\textcolor{QuartoInternalColor1}{}\textcolor{QuartoInternalColor1}{getindex}\textcolor{QuartoInternalColor1}{}\textcolor{QuartoInternalColor1}{}\textcolor{QuartoInternalColor1}{(}\textcolor{QuartoInternalColor1}{}\textcolor{QuartoInternalColor1}{df}\textcolor{QuartoInternalColor1}{::}\textcolor{QuartoInternalColor1}{DataFrame, ::}\textcolor{QuartoInternalColor1}{typeof(!), }\textcolor{QuartoInternalColor1}{col\_ind}\textcolor{QuartoInternalColor1}{::}\textcolor{QuartoInternalColor1}{Symbol}\textcolor{QuartoInternalColor1}{}\textcolor{QuartoInternalColor1}{)}\textcolor{QuartoInternalColor1}{}
\textcolor{QuartoInternalColor1}{}\textcolor{QuartoInternalColor1}{   @}\textcolor{QuartoInternalColor1}{ }\textcolor{QuartoInternalColor2}{DataFrames}\textcolor{QuartoInternalColor1}{ }\textcolor{QuartoInternalColor1}{\textasciitilde{}/.julia/packages/DataFrames/b4w9K/src/dataframe/}\textcolor{QuartoInternalColor1}{}\textcolor{QuartoInternalColor1}{}\textcolor{QuartoInternalColor1}{dataframe.jl:557}\textcolor{QuartoInternalColor1}{}\textcolor{QuartoInternalColor1}{}
\textcolor{QuartoInternalColor1}{ [4] }\textcolor{QuartoInternalColor1}{}\textcolor{QuartoInternalColor1}{getproperty}\textcolor{QuartoInternalColor1}{}\textcolor{QuartoInternalColor1}{}\textcolor{QuartoInternalColor1}{(}\textcolor{QuartoInternalColor1}{}\textcolor{QuartoInternalColor1}{df}\textcolor{QuartoInternalColor1}{::}\textcolor{QuartoInternalColor1}{DataFrame, }\textcolor{QuartoInternalColor1}{col\_ind}\textcolor{QuartoInternalColor1}{::}\textcolor{QuartoInternalColor1}{Symbol}\textcolor{QuartoInternalColor1}{}\textcolor{QuartoInternalColor1}{)}\textcolor{QuartoInternalColor1}{}
\textcolor{QuartoInternalColor1}{}\textcolor{QuartoInternalColor1}{   @}\textcolor{QuartoInternalColor1}{ }\textcolor{QuartoInternalColor2}{DataFrames}\textcolor{QuartoInternalColor1}{ }\textcolor{QuartoInternalColor1}{\textasciitilde{}/.julia/packages/DataFrames/b4w9K/src/abstractdataframe/}\textcolor{QuartoInternalColor1}{}\textcolor{QuartoInternalColor1}{}\textcolor{QuartoInternalColor1}{abstractdataframe.jl:448}\textcolor{QuartoInternalColor1}{}\textcolor{QuartoInternalColor1}{}
\textcolor{QuartoInternalColor1}{ [5] top-level scope}
\textcolor{QuartoInternalColor1}{}\textcolor{QuartoInternalColor1}{   @}\textcolor{QuartoInternalColor1}{ }\textcolor{QuartoInternalColor1}{\textasciitilde{}/Teaching/BEE4850/solutions/hw04/}\textcolor{QuartoInternalColor1}{}\textcolor{QuartoInternalColor1}{}\textcolor{QuartoInternalColor1}{hw04.qmd:231}\textcolor{QuartoInternalColor1}{}\textcolor{QuartoInternalColor1}{}
\end{Highlighting}

}

\caption{\label{fig-death-data-obs}}

\end{figure}%

We could fit the regression models in a few different ways (in Julia, we
could also use the \texttt{GLM.jl} package), but for consistency's sake
we'll do this with numerical optimization. First, let's fit a model for
the impact of time.

\phantomsection\label{annotated-cell-12}%
\begin{Shaded}
\begin{Highlighting}[]
\CommentTok{\# general log{-}likelihood function for a linear regression with covariate x}
\KeywordTok{function} \FunctionTok{timereg\_loglik}\NormalTok{(p, y, t)}
\NormalTok{    μ₀, μ₁, σ }\OperatorTok{=}\NormalTok{ p}
\NormalTok{    μ }\OperatorTok{=}\NormalTok{ μ₀ }\OperatorTok{.+}\NormalTok{ μ₁ }\OperatorTok{*}\NormalTok{ t}
\NormalTok{    ll }\OperatorTok{=} \FunctionTok{sum}\NormalTok{(}\FunctionTok{logpdf}\NormalTok{.(}\FunctionTok{Normal}\NormalTok{.(μ, σ), y))}
    \ControlFlowTok{return}\NormalTok{ ll}
\KeywordTok{end}

\NormalTok{lb }\OperatorTok{=}\NormalTok{ [}\FloatTok{0.0}\NormalTok{, }\OperatorTok{{-}}\FloatTok{100.0}\NormalTok{, }\FloatTok{0.1}\NormalTok{]}
\NormalTok{ub }\OperatorTok{=}\NormalTok{ [}\FloatTok{200.0}\NormalTok{, }\FloatTok{100.0}\NormalTok{, }\FloatTok{20.0}\NormalTok{]}
\NormalTok{x₀ }\OperatorTok{=}\NormalTok{ [}\FloatTok{100.0}\NormalTok{, }\FloatTok{0.0}\NormalTok{, }\FloatTok{5.0}\NormalTok{]}

\NormalTok{optim\_out }\OperatorTok{=}\NormalTok{ Optim.}\FunctionTok{optimize}\NormalTok{(p }\OperatorTok{{-}\textgreater{}} \FunctionTok{{-}timereg\_loglik}\NormalTok{(p, chicago\_dat.death, chicago\_dat.time), lb, ub, x₀)}
\NormalTok{ll\_time }\OperatorTok{=} \OperatorTok{{-}}\NormalTok{optim\_out.minimum }\hspace*{\fill}\NormalTok{\circled{1}}
\NormalTok{θ\_time }\OperatorTok{=}\NormalTok{ optim\_out.minimizer}
\PreprocessorTok{@show} \FunctionTok{round}\NormalTok{.(θ\_time; digits}\OperatorTok{=}\FloatTok{1}\NormalTok{);}
\end{Highlighting}
\end{Shaded}

\begin{description}
\tightlist
\item[\circled{1}]
Since we minimized the \emph{negative} log-likelihood, to get the
log-likelihood we need to take the negative of the minimum value again.
\end{description}

\begin{verbatim}
round.(θ_time; digits = 1) = [101.7, 0.0, 14.1]
\end{verbatim}

Plotting this regression over the data:

\begin{figure}

\centering{

\begin{Shaded}
\begin{Highlighting}[]
\FunctionTok{plot!}\NormalTok{(p\_deaths, chicago\_dat.}\DataTypeTok{Date}\NormalTok{, θ\_time[}\FloatTok{1}\NormalTok{] }\OperatorTok{.+}\NormalTok{ θ\_time[}\FloatTok{2}\NormalTok{] }\OperatorTok{*}\NormalTok{ chicago\_dat.time, lw}\OperatorTok{=}\FloatTok{2}\NormalTok{, linestyle}\OperatorTok{=:}\NormalTok{dash, color}\OperatorTok{=:}\NormalTok{red, label}\OperatorTok{=}\StringTok{"Time Regression"}\NormalTok{)}
\end{Highlighting}
\end{Shaded}

\begin{Highlighting}
\textcolor{black}{UndefVarError: UndefVarError(:p\_deaths, Main.Notebook)}
\textcolor{black}{UndefVarError: `p\_deaths` not defined in `Main.Notebook`}
\textcolor{black}{Suggestion: check for spelling errors or missing imports.}
\textcolor{black}{Stacktrace:}
\textcolor{black}{ [1] top-level scope}
\textcolor{black}{}\textcolor{QuartoInternalColor1}{   @}\textcolor{QuartoInternalColor1}{ }\textcolor{QuartoInternalColor1}{\textasciitilde{}/Teaching/BEE4850/solutions/hw04/}\textcolor{QuartoInternalColor1}{}\textcolor{QuartoInternalColor1}{}\textcolor{QuartoInternalColor1}{hw04.qmd:268}\textcolor{QuartoInternalColor1}{}\textcolor{QuartoInternalColor1}{}
\end{Highlighting}

}

\caption{\label{fig-time-regression}}

\end{figure}%

We can see that even though the time coefficient is close to zero,
suggesting a minimal time-dependent trend, over the course of the
dataset there is a slight increase over this time (seen in
Figure~\ref{fig-time-regression}).

\paragraph{Problem 2.2}\label{problem-2.2}

Now let's look at the impact of temperature.

\begin{Shaded}
\begin{Highlighting}[]
\CommentTok{\# general log{-}likelihood function for a linear regression with covariate x}
\KeywordTok{function} \FunctionTok{tempreg\_loglik}\NormalTok{(p, y, temp)}
\NormalTok{    μ₀, μ₁, σ }\OperatorTok{=}\NormalTok{ p}
\NormalTok{    μ }\OperatorTok{=}\NormalTok{ μ₀ }\OperatorTok{.+}\NormalTok{ μ₁ }\OperatorTok{*}\NormalTok{ temp}
\NormalTok{    ll }\OperatorTok{=} \FunctionTok{sum}\NormalTok{(}\FunctionTok{logpdf}\NormalTok{.(}\FunctionTok{Normal}\NormalTok{.(μ, σ), y))}
    \ControlFlowTok{return}\NormalTok{ ll}
\KeywordTok{end}

\NormalTok{lb }\OperatorTok{=}\NormalTok{ [}\FloatTok{0.0}\NormalTok{, }\OperatorTok{{-}}\FloatTok{100.0}\NormalTok{, }\FloatTok{0.1}\NormalTok{]}
\NormalTok{ub }\OperatorTok{=}\NormalTok{ [}\FloatTok{200.0}\NormalTok{, }\FloatTok{100.0}\NormalTok{, }\FloatTok{20.0}\NormalTok{]}
\NormalTok{x₀ }\OperatorTok{=}\NormalTok{ [}\FloatTok{100.0}\NormalTok{, }\FloatTok{0.0}\NormalTok{, }\FloatTok{5.0}\NormalTok{]}

\NormalTok{optim\_out }\OperatorTok{=}\NormalTok{ Optim.}\FunctionTok{optimize}\NormalTok{(p }\OperatorTok{{-}\textgreater{}} \FunctionTok{{-}tempreg\_loglik}\NormalTok{(p, chicago\_dat.death, chicago\_dat.tmpd), lb, ub, x₀)}
\NormalTok{ll\_temp }\OperatorTok{=} \OperatorTok{{-}}\NormalTok{optim\_out.minimum}
\NormalTok{θ\_temp }\OperatorTok{=}\NormalTok{ optim\_out.minimizer}
\PreprocessorTok{@show} \FunctionTok{round}\NormalTok{.(θ\_temp; digits}\OperatorTok{=}\FloatTok{1}\NormalTok{);}
\end{Highlighting}
\end{Shaded}

\begin{verbatim}
round.(θ_temp; digits = 1) = [129.2, -0.4, 12.4]
\end{verbatim}

Adding this regression line to the plot:

\begin{figure}

\centering{

\begin{Shaded}
\begin{Highlighting}[]
\FunctionTok{plot!}\NormalTok{(p\_deaths, chicago\_dat.}\DataTypeTok{Date}\NormalTok{, θ\_temp[}\FloatTok{1}\NormalTok{] }\OperatorTok{.+}\NormalTok{ θ\_temp[}\FloatTok{2}\NormalTok{] }\OperatorTok{*}\NormalTok{ chicago\_dat.tmpd, lw}\OperatorTok{=}\FloatTok{2}\NormalTok{, linestyle}\OperatorTok{=:}\NormalTok{dash, color}\OperatorTok{=:}\NormalTok{blue, label}\OperatorTok{=}\StringTok{"Temperature Regression"}\NormalTok{)}
\end{Highlighting}
\end{Shaded}

\begin{Highlighting}
\textcolor{black}{UndefVarError: UndefVarError(:p\_deaths, Main.Notebook)}
\textcolor{black}{UndefVarError: `p\_deaths` not defined in `Main.Notebook`}
\textcolor{black}{Suggestion: check for spelling errors or missing imports.}
\textcolor{black}{Stacktrace:}
\textcolor{black}{ [1] top-level scope}
\textcolor{black}{}\textcolor{QuartoInternalColor1}{   @}\textcolor{QuartoInternalColor1}{ }\textcolor{QuartoInternalColor1}{\textasciitilde{}/Teaching/BEE4850/solutions/hw04/}\textcolor{QuartoInternalColor1}{}\textcolor{QuartoInternalColor1}{}\textcolor{QuartoInternalColor1}{hw04.qmd:308}\textcolor{QuartoInternalColor1}{}\textcolor{QuartoInternalColor1}{}
\end{Highlighting}

}

\caption{\label{fig-temp-regression}}

\end{figure}%

Figure~\ref{fig-temp-regression} shows that the temperature regression
results in some oscillations due to the seasonal impact of temperature,
which visually appear to align with the oscillations in the deaths data.

\paragraph{Problem 2.3}\label{problem-2.3}

Our last model adds PM2.5 concentrations to the previous temperature
regression. Here we need to drop the entries corresponding to missing
data.

\begin{Shaded}
\begin{Highlighting}[]
\CommentTok{\# general log{-}likelihood function for a linear regression with covariate x}
\KeywordTok{function} \FunctionTok{multireg\_loglik}\NormalTok{(p, y, temp, pm25)}
\NormalTok{    μ₀, μ₁, μ₂, σ }\OperatorTok{=}\NormalTok{ p}
\NormalTok{    μ }\OperatorTok{=}\NormalTok{ μ₀ }\OperatorTok{.+}\NormalTok{ μ₁ }\OperatorTok{*}\NormalTok{ temp }\OperatorTok{.+}\NormalTok{ μ₂ }\OperatorTok{*}\NormalTok{ pm25}
\NormalTok{    ll }\OperatorTok{=} \FunctionTok{sum}\NormalTok{(}\FunctionTok{logpdf}\NormalTok{.(}\FunctionTok{Normal}\NormalTok{.(μ, σ), y))}
    \ControlFlowTok{return}\NormalTok{ ll}
\KeywordTok{end}

\NormalTok{lb }\OperatorTok{=}\NormalTok{ [}\FloatTok{0.0}\NormalTok{, }\OperatorTok{{-}}\FloatTok{100.0}\NormalTok{, }\OperatorTok{{-}}\FloatTok{100.0}\NormalTok{, }\FloatTok{0.1}\NormalTok{]}
\NormalTok{ub }\OperatorTok{=}\NormalTok{ [}\FloatTok{200.0}\NormalTok{, }\FloatTok{100.0}\NormalTok{, }\FloatTok{100.0}\NormalTok{, }\FloatTok{20.0}\NormalTok{]}
\NormalTok{x₀ }\OperatorTok{=}\NormalTok{ [}\FloatTok{100.0}\NormalTok{, }\FloatTok{0.0}\NormalTok{, }\FloatTok{0.0}\NormalTok{, }\FloatTok{5.0}\NormalTok{]}

\NormalTok{optim\_out }\OperatorTok{=}\NormalTok{ Optim.}\FunctionTok{optimize}\NormalTok{(p }\OperatorTok{{-}\textgreater{}} \FunctionTok{{-}multireg\_loglik}\NormalTok{(p, chicago\_dat.death, chicago\_dat.tmpd, chicago\_dat.pm25median), lb, ub, x₀)}
\NormalTok{ll\_multi }\OperatorTok{=} \OperatorTok{{-}}\NormalTok{optim\_out.minimum}
\NormalTok{θ\_multi }\OperatorTok{=}\NormalTok{ optim\_out.minimizer}
\PreprocessorTok{@show} \FunctionTok{round}\NormalTok{.(θ\_multi; digits}\OperatorTok{=}\FloatTok{1}\NormalTok{);}
\end{Highlighting}
\end{Shaded}

\begin{verbatim}
round.(θ_multi; digits = 1) = [129.3, -0.4, 0.2, 12.3]
\end{verbatim}

We add this regression to the plot:

\begin{figure}

\centering{

\begin{Shaded}
\begin{Highlighting}[]
\FunctionTok{plot!}\NormalTok{(p\_deaths, chicago\_dat.}\DataTypeTok{Date}\NormalTok{, θ\_multi[}\FloatTok{1}\NormalTok{] }\OperatorTok{.+}\NormalTok{ θ\_multi[}\FloatTok{2}\NormalTok{] }\OperatorTok{*}\NormalTok{ chicago\_dat.tmpd }\OperatorTok{.+}\NormalTok{ θ\_multi[}\FloatTok{3}\NormalTok{] }\OperatorTok{*}\NormalTok{ chicago\_dat.pm25median, lw}\OperatorTok{=}\FloatTok{2}\NormalTok{, linestyle}\OperatorTok{=:}\NormalTok{dash, color}\OperatorTok{=:}\NormalTok{violet, label}\OperatorTok{=}\StringTok{"Temperature and PM2.5 Regression"}\NormalTok{)}
\end{Highlighting}
\end{Shaded}

\begin{Highlighting}
\textcolor{black}{UndefVarError: UndefVarError(:p\_deaths, Main.Notebook)}
\textcolor{black}{UndefVarError: `p\_deaths` not defined in `Main.Notebook`}
\textcolor{black}{Suggestion: check for spelling errors or missing imports.}
\textcolor{black}{Stacktrace:}
\textcolor{black}{ [1] top-level scope}
\textcolor{black}{}\textcolor{QuartoInternalColor1}{   @}\textcolor{QuartoInternalColor1}{ }\textcolor{QuartoInternalColor1}{\textasciitilde{}/Teaching/BEE4850/solutions/hw04/}\textcolor{QuartoInternalColor1}{}\textcolor{QuartoInternalColor1}{}\textcolor{QuartoInternalColor1}{hw04.qmd:347}\textcolor{QuartoInternalColor1}{}\textcolor{QuartoInternalColor1}{}
\end{Highlighting}

}

\caption{\label{fig-multi-regression}}

\end{figure}%

\paragraph{Problem 2.4}\label{problem-2.4}

Now we want to conduct a 5-fold cross-validation to compare each model's
performance. As set up, we don't need to account for any particular
structure, as we aren't accounting for any underlying seasonal effects
that aren't captured by the covariates.

We will write a function to conduct these for each of our models.

\phantomsection\label{annotated-cell-18}%
\begin{Shaded}
\begin{Highlighting}[]
\KeywordTok{function} \FunctionTok{time\_cv}\NormalTok{(dat, nfolds}\OperatorTok{=}\FloatTok{5}\NormalTok{)}
    \CommentTok{\# split the data}
\NormalTok{    labels }\OperatorTok{=} \FunctionTok{repeat}\NormalTok{(}\FloatTok{1}\OperatorTok{:}\NormalTok{nfolds, outer }\OperatorTok{=} \FunctionTok{Int}\NormalTok{(}\FunctionTok{floor}\NormalTok{(}\FunctionTok{nrow}\NormalTok{(dat) }\OperatorTok{/}\NormalTok{ nfolds))) }\hspace*{\fill}\NormalTok{\circled{1}}
\NormalTok{    labels }\OperatorTok{=}\NormalTok{ [labels; }\FloatTok{1}\NormalTok{; }\FloatTok{2}\NormalTok{] }
\NormalTok{    fold\_labels }\OperatorTok{=} \FunctionTok{shuffle}\NormalTok{(labels) }\CommentTok{\# randomly permite the labels}
    \CommentTok{\# create storage for error metrics}
\NormalTok{    cv\_ll }\OperatorTok{=} \FunctionTok{zeros}\NormalTok{(nfolds)}
\NormalTok{    cv\_mse }\OperatorTok{=} \FunctionTok{zeros}\NormalTok{(nfolds)}
    \CommentTok{\# set optimization parameters}
\NormalTok{    lb }\OperatorTok{=}\NormalTok{ [}\FloatTok{0.0}\NormalTok{, }\OperatorTok{{-}}\FloatTok{100.0}\NormalTok{, }\FloatTok{0.1}\NormalTok{]}
\NormalTok{    ub }\OperatorTok{=}\NormalTok{ [}\FloatTok{200.0}\NormalTok{, }\FloatTok{100.0}\NormalTok{, }\FloatTok{20.0}\NormalTok{]}
\NormalTok{    p0 }\OperatorTok{=}\NormalTok{ [}\FloatTok{100.0}\NormalTok{, }\FloatTok{0.0}\NormalTok{, }\FloatTok{5.0}\NormalTok{]}
    \ControlFlowTok{for}\NormalTok{ fold }\KeywordTok{in} \FloatTok{1}\OperatorTok{:}\NormalTok{nfolds}
\NormalTok{        test\_idx }\OperatorTok{=} \FunctionTok{findall}\NormalTok{(fold\_labels }\OperatorTok{.==}\NormalTok{ fold)}
\NormalTok{        train\_idx }\OperatorTok{=} \FunctionTok{setdiff}\NormalTok{(}\FloatTok{1}\OperatorTok{:}\FunctionTok{nrow}\NormalTok{(dat), test\_idx)}
        \CommentTok{\# fit model to the fold}
\NormalTok{        optim\_out }\OperatorTok{=} \FunctionTok{optimize}\NormalTok{(p }\OperatorTok{{-}\textgreater{}} \FunctionTok{{-}timereg\_loglik}\NormalTok{(p, dat[train\_idx, }\OperatorTok{:}\NormalTok{death], dat[train\_idx, }\OperatorTok{:}\NormalTok{time]), lb, ub, p0)}
\NormalTok{        p\_mle }\OperatorTok{=}\NormalTok{ optim\_out.minimizer}
\NormalTok{        pred }\OperatorTok{=}\NormalTok{ p\_mle[}\FloatTok{1}\NormalTok{] }\OperatorTok{.+}\NormalTok{ p\_mle[}\FloatTok{2}\NormalTok{] }\OperatorTok{*}\NormalTok{ dat[test\_idx, }\OperatorTok{:}\NormalTok{time]}
\NormalTok{        cv\_ll[fold] }\OperatorTok{=} \FunctionTok{{-}mean}\NormalTok{(}\FunctionTok{logpdf}\NormalTok{.(}\FunctionTok{Normal}\NormalTok{.(pred, p\_mle[}\FloatTok{3}\NormalTok{]), dat[test\_idx, }\OperatorTok{:}\NormalTok{death])) }\hspace*{\fill}\NormalTok{\circled{2}}
\NormalTok{        cv\_mse[fold] }\OperatorTok{=} \FunctionTok{mean}\NormalTok{((pred }\OperatorTok{.{-}}\NormalTok{ dat[test\_idx, }\OperatorTok{:}\NormalTok{death])}\OperatorTok{.\^{}}\FloatTok{2}\NormalTok{)}
    \ControlFlowTok{end}
    \ControlFlowTok{return}\NormalTok{ cv\_ll, cv\_mse}
\KeywordTok{end}

\KeywordTok{function} \FunctionTok{temp\_cv}\NormalTok{(dat, nfolds}\OperatorTok{=}\FloatTok{5}\NormalTok{)}
    \CommentTok{\# split the data}
\NormalTok{    labels }\OperatorTok{=} \FunctionTok{repeat}\NormalTok{(}\FloatTok{1}\OperatorTok{:}\NormalTok{nfolds, outer }\OperatorTok{=} \FunctionTok{Int}\NormalTok{(}\FunctionTok{floor}\NormalTok{(}\FunctionTok{nrow}\NormalTok{(dat) }\OperatorTok{/}\NormalTok{ nfolds))) }\hspace*{\fill}\NormalTok{\circled{1}}
\NormalTok{    labels }\OperatorTok{=}\NormalTok{ [labels; }\FloatTok{1}\NormalTok{; }\FloatTok{2}\NormalTok{] }
\NormalTok{    fold\_labels }\OperatorTok{=} \FunctionTok{shuffle}\NormalTok{(labels) }\CommentTok{\# randomly permite the labels}
    \CommentTok{\# create storage for error metrics}
\NormalTok{    cv\_ll }\OperatorTok{=} \FunctionTok{zeros}\NormalTok{(nfolds)}
\NormalTok{    cv\_mse }\OperatorTok{=} \FunctionTok{zeros}\NormalTok{(nfolds)}
    \CommentTok{\# set optimization parameters}
\NormalTok{    lb }\OperatorTok{=}\NormalTok{ [}\FloatTok{0.0}\NormalTok{, }\OperatorTok{{-}}\FloatTok{100.0}\NormalTok{, }\FloatTok{0.1}\NormalTok{]}
\NormalTok{    ub }\OperatorTok{=}\NormalTok{ [}\FloatTok{200.0}\NormalTok{, }\FloatTok{100.0}\NormalTok{, }\FloatTok{20.0}\NormalTok{]}
\NormalTok{    p0 }\OperatorTok{=}\NormalTok{ [}\FloatTok{100.0}\NormalTok{, }\FloatTok{0.0}\NormalTok{, }\FloatTok{5.0}\NormalTok{]}
    \ControlFlowTok{for}\NormalTok{ fold }\KeywordTok{in} \FloatTok{1}\OperatorTok{:}\NormalTok{nfolds}
\NormalTok{        test\_idx }\OperatorTok{=} \FunctionTok{findall}\NormalTok{(fold\_labels }\OperatorTok{.==}\NormalTok{ fold)}
\NormalTok{        train\_idx }\OperatorTok{=} \FunctionTok{setdiff}\NormalTok{(}\FloatTok{1}\OperatorTok{:}\FunctionTok{nrow}\NormalTok{(dat), test\_idx)}
        \CommentTok{\# fit model to the fold}
\NormalTok{        optim\_out }\OperatorTok{=} \FunctionTok{optimize}\NormalTok{(p }\OperatorTok{{-}\textgreater{}} \FunctionTok{{-}tempreg\_loglik}\NormalTok{(p, dat[train\_idx, }\OperatorTok{:}\NormalTok{death], dat[train\_idx, }\OperatorTok{:}\NormalTok{tmpd]), lb, ub, p0)}
\NormalTok{        p\_mle }\OperatorTok{=}\NormalTok{ optim\_out.minimizer}
\NormalTok{        pred }\OperatorTok{=}\NormalTok{ p\_mle[}\FloatTok{1}\NormalTok{] }\OperatorTok{.+}\NormalTok{ p\_mle[}\FloatTok{2}\NormalTok{] }\OperatorTok{*}\NormalTok{ dat[test\_idx, }\OperatorTok{:}\NormalTok{tmpd]}
\NormalTok{        cv\_ll[fold] }\OperatorTok{=} \FunctionTok{{-}mean}\NormalTok{(}\FunctionTok{logpdf}\NormalTok{.(}\FunctionTok{Normal}\NormalTok{.(pred, p\_mle[}\FloatTok{3}\NormalTok{]), dat[test\_idx, }\OperatorTok{:}\NormalTok{death])) }\hspace*{\fill}\NormalTok{\circled{2}}
\NormalTok{        cv\_mse[fold] }\OperatorTok{=} \FunctionTok{mean}\NormalTok{((pred }\OperatorTok{.{-}}\NormalTok{ dat[test\_idx, }\OperatorTok{:}\NormalTok{death])}\OperatorTok{.\^{}}\FloatTok{2}\NormalTok{)}
    \ControlFlowTok{end}
    \ControlFlowTok{return}\NormalTok{ cv\_ll, cv\_mse}
\KeywordTok{end}

\KeywordTok{function} \FunctionTok{temp\_pm\_cv}\NormalTok{(dat, nfolds}\OperatorTok{=}\FloatTok{5}\NormalTok{)}
    \CommentTok{\# split the data}
\NormalTok{    labels }\OperatorTok{=} \FunctionTok{repeat}\NormalTok{(}\FloatTok{1}\OperatorTok{:}\NormalTok{nfolds, outer }\OperatorTok{=} \FunctionTok{Int}\NormalTok{(}\FunctionTok{floor}\NormalTok{(}\FunctionTok{nrow}\NormalTok{(dat) }\OperatorTok{/}\NormalTok{ nfolds))) }\hspace*{\fill}\NormalTok{\circled{1}}
\NormalTok{    labels }\OperatorTok{=}\NormalTok{ [labels; }\FloatTok{1}\NormalTok{; }\FloatTok{2}\NormalTok{] }
\NormalTok{    fold\_labels }\OperatorTok{=} \FunctionTok{shuffle}\NormalTok{(labels) }\CommentTok{\# randomly permite the labels}
    \CommentTok{\# create storage for error metrics}
\NormalTok{    cv\_ll }\OperatorTok{=} \FunctionTok{zeros}\NormalTok{(nfolds)}
\NormalTok{    cv\_mse }\OperatorTok{=} \FunctionTok{zeros}\NormalTok{(nfolds)}
    \CommentTok{\# set optimization parameters}
\NormalTok{    lb }\OperatorTok{=}\NormalTok{ [}\FloatTok{0.0}\NormalTok{, }\OperatorTok{{-}}\FloatTok{100.0}\NormalTok{, }\OperatorTok{{-}}\FloatTok{100.0}\NormalTok{, }\FloatTok{0.1}\NormalTok{]}
\NormalTok{    ub }\OperatorTok{=}\NormalTok{ [}\FloatTok{200.0}\NormalTok{, }\FloatTok{100.0}\NormalTok{, }\FloatTok{100.0}\NormalTok{, }\FloatTok{20.0}\NormalTok{]}
\NormalTok{    p0 }\OperatorTok{=}\NormalTok{ [}\FloatTok{100.0}\NormalTok{, }\FloatTok{0.0}\NormalTok{, }\FloatTok{0.0}\NormalTok{, }\FloatTok{5.0}\NormalTok{]}

    \ControlFlowTok{for}\NormalTok{ fold }\KeywordTok{in} \FloatTok{1}\OperatorTok{:}\NormalTok{nfolds}
\NormalTok{        test\_idx }\OperatorTok{=} \FunctionTok{findall}\NormalTok{(fold\_labels }\OperatorTok{.==}\NormalTok{ fold)}
\NormalTok{        train\_idx }\OperatorTok{=} \FunctionTok{setdiff}\NormalTok{(}\FloatTok{1}\OperatorTok{:}\FunctionTok{nrow}\NormalTok{(dat), test\_idx)}
        \CommentTok{\# fit model to the fold}
\NormalTok{        optim\_out }\OperatorTok{=} \FunctionTok{optimize}\NormalTok{(p }\OperatorTok{{-}\textgreater{}} \FunctionTok{{-}multireg\_loglik}\NormalTok{(p, dat[train\_idx, }\OperatorTok{:}\NormalTok{death], dat[train\_idx, }\OperatorTok{:}\NormalTok{tmpd], dat[train\_idx, }\OperatorTok{:}\NormalTok{pm25median]), lb, ub, p0)}
\NormalTok{        p\_mle }\OperatorTok{=}\NormalTok{ optim\_out.minimizer}
\NormalTok{        pred }\OperatorTok{=}\NormalTok{ p\_mle[}\FloatTok{1}\NormalTok{] }\OperatorTok{.+}\NormalTok{ p\_mle[}\FloatTok{2}\NormalTok{] }\OperatorTok{*}\NormalTok{ dat[test\_idx, }\OperatorTok{:}\NormalTok{tmpd] }\OperatorTok{.+}\NormalTok{ p\_mle[}\FloatTok{3}\NormalTok{] }\OperatorTok{*}\NormalTok{ dat[test\_idx, }\OperatorTok{:}\NormalTok{pm25median]}
\NormalTok{        cv\_ll[fold] }\OperatorTok{=} \FunctionTok{{-}mean}\NormalTok{(}\FunctionTok{logpdf}\NormalTok{.(}\FunctionTok{Normal}\NormalTok{.(pred, p\_mle[}\FloatTok{4}\NormalTok{]), dat[test\_idx, }\OperatorTok{:}\NormalTok{death])) }\hspace*{\fill}\NormalTok{\circled{2}}
\NormalTok{        cv\_mse[fold] }\OperatorTok{=} \FunctionTok{mean}\NormalTok{((pred }\OperatorTok{.{-}}\NormalTok{ dat[test\_idx, }\OperatorTok{:}\NormalTok{death])}\OperatorTok{.\^{}}\FloatTok{2}\NormalTok{)}
    \ControlFlowTok{end}
    \ControlFlowTok{return}\NormalTok{ cv\_ll, cv\_mse}

\KeywordTok{end}
\end{Highlighting}
\end{Shaded}

\begin{verbatim}
temp_pm_cv (generic function with 2 methods)
\end{verbatim}

Now let's compare the results for each model.

\begin{Shaded}
\begin{Highlighting}[]
\NormalTok{time\_cv\_ll, time\_cv\_mse }\OperatorTok{=} \FunctionTok{time\_cv}\NormalTok{(chicago\_dat, }\FloatTok{5}\NormalTok{)}
\NormalTok{temp\_cv\_ll, temp\_cv\_mse }\OperatorTok{=} \FunctionTok{temp\_cv}\NormalTok{(chicago\_dat, }\FloatTok{5}\NormalTok{)}
\NormalTok{multi\_cv\_ll, multi\_cv\_mse }\OperatorTok{=} \FunctionTok{temp\_pm\_cv}\NormalTok{(chicago\_dat, }\FloatTok{5}\NormalTok{)}

\NormalTok{cv\_models }\OperatorTok{=}\NormalTok{ [}\StringTok{"Time"}\NormalTok{, }\StringTok{"Temperature"}\NormalTok{, }\StringTok{"Temperature + PM 2.5"}\NormalTok{]}
\NormalTok{cv\_ll\_all }\OperatorTok{=}\NormalTok{ [}\FunctionTok{mean}\NormalTok{(time\_cv\_ll), }\FunctionTok{mean}\NormalTok{(temp\_cv\_ll), }\FunctionTok{mean}\NormalTok{(multi\_cv\_ll)]}
\NormalTok{cv\_mse\_all }\OperatorTok{=}\NormalTok{ [}\FunctionTok{mean}\NormalTok{(time\_cv\_mse), }\FunctionTok{mean}\NormalTok{(temp\_cv\_mse), }\FunctionTok{mean}\NormalTok{(multi\_cv\_mse)]}

\NormalTok{cv\_results }\OperatorTok{=} \FunctionTok{DataFrame}\NormalTok{(Model}\OperatorTok{=}\NormalTok{cv\_models, LogScore}\OperatorTok{=}\NormalTok{cv\_ll\_all, MSE}\OperatorTok{=}\NormalTok{cv\_mse\_all)}
\end{Highlighting}
\end{Shaded}

\begin{tabular}{r|ccc}
    & Model & LogScore & MSE\\
    \hline
    & String & Float64 & Float64\\
    \hline
    1 & Time & 4.07372 & 201.236 \\
    2 & Temperature & 3.94227 & 155.441 \\
    3 & Temperature + PM 2.5 & 3.93076 & 151.581 \\
\end{tabular}

By these scores, the temperature + PM 2.5 model scores best for both the
log score and the MSE, so we might conclude that it would generalize the
best out-of-sample. The temperature-only model performs similarly, and
the time model, as we saw in Problem 2.1, does not perform particularly
well.

\paragraph{Problem 2.5}\label{problem-2.5}

To compute the AIC values we will use the log-likelihoods we computed
when we fit each model.

\begin{Shaded}
\begin{Highlighting}[]
\NormalTok{aic\_time }\OperatorTok{=} \OperatorTok{{-}}\FloatTok{2} \OperatorTok{*}\NormalTok{ (ll\_time }\OperatorTok{{-}} \FloatTok{3}\NormalTok{)}
\NormalTok{aic\_temp }\OperatorTok{=} \OperatorTok{{-}}\FloatTok{2} \OperatorTok{*}\NormalTok{ (ll\_temp }\OperatorTok{{-}} \FloatTok{3}\NormalTok{)}
\NormalTok{aic\_multi }\OperatorTok{=} \OperatorTok{{-}}\FloatTok{2} \OperatorTok{*}\NormalTok{ (ll\_multi }\OperatorTok{{-}} \FloatTok{4}\NormalTok{)}
\NormalTok{aic\_df }\OperatorTok{=} \FunctionTok{DataFrame}\NormalTok{(Models}\OperatorTok{=}\NormalTok{[}\StringTok{"Time"}\NormalTok{, }\StringTok{"Temperature"}\NormalTok{, }\StringTok{"Temperature and PM2.5"}\NormalTok{], AIC}\OperatorTok{=}\FunctionTok{round}\NormalTok{.([aic\_time, aic\_temp, aic\_multi]; digits}\OperatorTok{=}\FloatTok{0}\NormalTok{))}
\end{Highlighting}
\end{Shaded}

\begin{tabular}{r|cc}
    & Models & AIC\\
    \hline
    & String & Float64\\
    \hline
    1 & Time & 5918.0 \\
    2 & Temperature & 5729.0 \\
    3 & Temperature and PM2.5 & 5714.0 \\
\end{tabular}

We can see that the Temperature and PM2.5 model minimizes AIC, and the
\(\Delta \text{AIC}\) for the temperature-only model is around 15. These
differences suggest that there is relatively strong evidence for the
influence of both temperature and PM2.5 compared to a temperature-only
model, as the predictive skill outweighs the potential overfitting from
adding the PM2.5 term. This is similar in direction to what we saw from
the cross-validation as well.

When comparing this quantitative result to the plots, drawing this
conclusion from Figure~\ref{fig-multi-regression} itself is challenging,
even if we might be able to squint and see that the fluctuations from
the PM2.5 model align slightly better with the temperature-only model.




\end{document}
